\section{Dependent types, with examples}\label{sec:01-basics:03-dependent-types:1}

Every function seen so far has a \emph{fixed} codomain: \texttt{double : Nat → Nat} returns a \texttt{Nat} regardless of the input \texttt{n}. Lean's type theory allows something more general: a type that itself depends on a \emph{value}, and a function whose return type changes depending on which argument it was given. This is a \textbf{dependent type}. It is the single feature that separates Lean (and other proof assistants) from an ordinary typed programming language, so it is worth building up slowly, from the smallest possible example, before naming the general pattern.

\subsection{The problem, in Python first}\label{sec:01-basics:03-dependent-types:2}

Before any Lean syntax, here is the actual problem dependent types solve, in a language with no type checker enforcing anything at all. Suppose a dot-product function is written over two lists:

\begin{lstlisting}[language=Python,style=python]
def dot(xs, ys):
    return sum(x * y for x, y in zip(xs, ys))

dot([1, 2, 3], [4, 5, 6])   # 32 (*\textemdash{}*) correct
dot([1, 2, 3], [4, 5])      # 14 (*\textemdash{}*) silently wrong: zip() truncates to the shorter list
\end{lstlisting}

This is worse than a crash. \texttt{dot([1, 2, 3], [4, 5])} does not raise anything --- \texttt{zip} just quietly drops the third element of \texttt{xs} and hands back a number that \emph{looks} like a perfectly good answer. The bug (calling \texttt{dot} on two lists of different lengths, which is mathematically nonsensical) is never caught, not by Python, not by a test that happens not to cover this call site, not by anything short of a human noticing the number looks off.

The underlying mistake is that a Python \texttt{list} carries no information about its own length in any way Python's function signatures can see or enforce. \texttt{xs : list} and \texttt{ys : list} say nothing about whether the two lists agree in size --- that fact exists only as a comment at best, never checked by anything. What is actually wanted is a signature saying ``\texttt{dot} accepts two lists \emph{of the same length} \texttt{n}, for \emph{any} \texttt{n}'' --- a signature that mentions and constrains a value (\texttt{n}), not just a type (\texttt{list}). Ordinary type systems, Python's included even with type hints added, have no way to say that. This is exactly the gap dependent types close, and the rest of this section builds the machinery to say it, starting from the smallest possible example.

\subsection{First: a type family, one type per number}\label{sec:01-basics:03-dependent-types:3}

Here is the smallest genuine example of a dependent type, and it is already sitting in Lean's core library. \texttt{Fin n} is the type of natural numbers \emph{strictly less than} \texttt{n}:

\begin{lstlisting}
#check Fin 3   -- Fin 3 : Type
#check Fin 5   -- Fin 5 : Type
\end{lstlisting}

\texttt{Fin} itself is not one type. It is a \textbf{recipe that produces a type once handed a number}: \texttt{Fin 3} and \texttt{Fin 5} are both genuine types, but they are \emph{different} types --- \texttt{Fin 3} has exactly three inhabitants (the numbers \texttt{0}, \texttt{1}, \texttt{2}) and \texttt{Fin 5} has exactly five. Compare this to something already known not to be dependent, \texttt{List α}: \texttt{List Nat} and \texttt{List Bool} are different types too, but only because \texttt{Nat} and \texttt{Bool} are different \emph{types} fed in for \texttt{α}. \texttt{Fin} is a different kind of thing. \texttt{Fin 3} and \texttt{Fin 5} differ even though \texttt{3} and \texttt{5} are both perfectly ordinary terms of the exact same type, \texttt{Nat}. The type \texttt{Fin} produces depends on which \emph{value} it is given, not just which type. That value-dependence is the entire definition of ``dependent type,'' and \texttt{Fin} gives it in the smallest possible dose.

The construction of \texttt{Fin} can be inspected directly:

\begin{lstlisting}
#print Fin
-- structure Fin (n : Nat) : Type
-- fields:
--   Fin.val  : Nat
--   Fin.isLt : val < n
\end{lstlisting}

So a term of \texttt{Fin n} is literally a pair: a \texttt{Nat} value, plus a \emph{proof} that the value is below \texttt{n}. The proof's very statement (\texttt{val < n}) mentions \texttt{n}, the value supplied. Change \texttt{n} and the result is a genuinely different type, with a genuinely different proof obligation attached. This bundling --- data, plus a proof whose statement depends on that data --- is the second half of the dependent-types story (formalized later in this section as a \textbf{Σ-type}); \texttt{Fin} is a real, live example of it, not a made-up one.

\subsection{\texorpdfstring{Second: a function whose \emph{return type} changes with its argument}{Second: a function whose return type changes with its argument}}\label{sec:01-basics:03-dependent-types:4}

Now for the companion idea: a \textbf{dependent function}, one whose return type depends on the specific \emph{value} of its argument, not just the argument's type. Define fixed-length vectors from scratch, the standard first example in any dependent-type-theory course:

\begin{lstlisting}
inductive Vec ((*$\alpha$*) : Type) : Nat (*$\rightarrow$*) Type where
  | nil  : Vec (*$\alpha$*) 0
  | cons : (*$\alpha$*) (*$\rightarrow$*) Vec (*$\alpha$*) n (*$\rightarrow$*) Vec (*$\alpha$*) (n + 1)
\end{lstlisting}

Read this exactly like \texttt{Nat}'s two-constructor definition from the previous section, with one new ingredient: \texttt{Vec α} is not a single type, it is a \textbf{family of types indexed by a \texttt{Nat}} --- \texttt{Vec α 0}, \texttt{Vec α 1}, \texttt{Vec α 2}, \ldots{} are all different types, one per length, and the length is tracked \emph{in the type itself}, not just at runtime. \texttt{nil} builds the unique length-\texttt{0} vector; \texttt{cons} takes an element and a length-\texttt{n} vector and produces a length-\texttt{(n+1)} vector --- the \texttt{n} used on both sides of \texttt{cons}'s arrow is the \emph{same} \texttt{n}, so the constructor itself enforces ``one longer than whatever it started with.''

Here is a function that \emph{builds} one of these, and its type is the dependent-function payoff:

\begin{lstlisting}
def Vec.replicate (a : (*$\alpha$*)) : (n : Nat) (*$\rightarrow$*) Vec (*$\alpha$*) n
  | 0     => Vec.nil
  | n + 1 => Vec.cons a (Vec.replicate a n)

#check @Vec.replicate
-- @Vec.replicate : {(*$\alpha$*) : Type} (*$\rightarrow$*) (*$\alpha$*) (*$\rightarrow$*) (n : Nat) (*$\rightarrow$*) Vec (*$\alpha$*) n

#eval (Vec.replicate 7 3 : Vec Nat 3)
-- Vec.cons 7 (Vec.cons 7 (Vec.cons 7 Vec.nil))
\end{lstlisting}

Look closely at the type \texttt{(n : Nat) → Vec α n}. The \texttt{n} that appears on the \emph{left} of the arrow (the argument) reappears inside the type on the \emph{right} of the arrow (the result). Feed \texttt{Vec.replicate a} the number \texttt{3}, and the result has type \texttt{Vec α 3}, specifically. Feed it \texttt{5}, and the very same function returns something of type \texttt{Vec α 5} instead. This is categorically different from \texttt{double : Nat → Nat}, where the output type (\texttt{Nat}) is fixed in advance and never reads the input value at all. Here, the \emph{type itself} changes based on which number was passed in. That is exactly what ``the codomain depends on the argument'' means, made as concrete as possible.

\subsection{Why bother: invariants become part of the type, not a side promise}\label{sec:01-basics:03-dependent-types:5}

The payoff is not just bookkeeping. Because the length lives in the type, Lean can rule out a whole class of mistakes \emph{before running anything at all}. Define a function that reads a vector's first element, which only makes sense for a \emph{non-empty} vector:

\begin{lstlisting}
def Vec.head : Vec (*$\alpha$*) (n + 1) (*$\rightarrow$*) (*$\alpha$*)
  | Vec.cons a _ => a
\end{lstlisting}

The argument type \texttt{Vec α (n + 1)} says, in the type itself, ``this only accepts vectors of length \emph{at least one}'' --- there is no separate runtime check for emptiness anywhere in this definition, because none is needed. Calling it on an empty vector is rejected before the expression ever runs:

\begin{lstlisting}
#check Vec.head Vec.nil
\end{lstlisting}

\begin{lstlisting}
error: Application type mismatch: The argument
  Vec.nil
has type
  Vec ?m 0
but is expected to have type
  Vec ?m (?n + 1)
in the application
  Vec.head Vec.nil
\end{lstlisting}

Nothing about ``index out of range'' happens at runtime, because the bad call is not a well-typed term in the first place --- the same ``ruled out before running'' guarantee from \hyperref[sec:01-basics:01-everything-has-a-type:1]{§1}, now enforced by an invariant (non-emptiness) that an ordinary, non-dependent type could not have expressed at all. \texttt{List α} has no way to say ``and this one is non-empty'' as part of its type; \texttt{Vec α (n+1)} says exactly that, for free, using only the machinery already on the table.

\textbf{Now close the loop on the Python example from the start of this section.} Here is \texttt{dot}, rewritten for \texttt{Vec} instead of Python's \texttt{list}, with both arguments required to share the \emph{same} length \texttt{n}:

\begin{lstlisting}
def Vec.dot : Vec Nat n (*$\rightarrow$*) Vec Nat n (*$\rightarrow$*) Nat
  | Vec.nil, Vec.nil => 0
  | Vec.cons x xs, Vec.cons y ys => x * y + Vec.dot xs ys
\end{lstlisting}

The signature \texttt{Vec Nat n → Vec Nat n → Nat} uses the \emph{same} \texttt{n} for both arguments --- that is not a naming coincidence, it is the whole point. Try to reproduce Python's silent bug:

\begin{lstlisting}
def v3 : Vec Nat 3 := Vec.cons 1 (Vec.cons 2 (Vec.cons 3 Vec.nil))
def v2 : Vec Nat 2 := Vec.cons 4 (Vec.cons 5 Vec.nil)

#check Vec.dot v3 v2
\end{lstlisting}

\begin{lstlisting}
error: Application type mismatch: The argument
  v2
has type
  Vec Nat 2
but is expected to have type
  Vec Nat 3
in the application
  v3.dot v2
\end{lstlisting}

Where Python's \texttt{dot([1,2,3], [4,5])} silently returned \texttt{14} --- a wrong answer with no error at all --- Lean's version does not even compile. The length-mismatch bug is not caught by a clever runtime check \emph{added} to \texttt{Vec.dot}; there is no such check anywhere in its three-line definition. It is caught because ``both arguments have the same length'' was stated once, in the type, and Lean enforces every type it is given, automatically, for every call site, without exception.

(Lean's actual, built-in \texttt{Vector α n} --- distinct from the toy \texttt{Vec} built here --- is defined differently under the hood, as an \texttt{Array α} paired with a proof that its size equals \texttt{n}, for performance reasons, the same way \texttt{Nat}'s \emph{presentation} above as Peano's \texttt{zero}/\texttt{succ} does not reflect how Lean actually stores numbers at runtime (as fast arbitrary-precision integers). The \texttt{Vec} built in this section is the traditional textbook definition --- simpler to reason about, and the one every type-theory reference uses first --- while Lean's real \texttt{Vector} is engineered for speed. Both are dependent types in exactly the sense described here.)

\subsection{The general pattern: Π-types}\label{sec:01-basics:03-dependent-types:6}

Both examples above are instances of one idea. A \textbf{dependent function type}, written with \(\Pi\) (``Pi,'' for ``dependent product''), generalizes the ordinary function-space \(A \to B\):

\[
\prod_{x : A} B(x)
\]

read: ``a function that, given any \(x : A\), returns a term of type \(B(x)\)'' --- a type allowed to mention \(x\). When \(B(x)\) does not actually depend on \(x\), this collapses exactly to the ordinary function type \(A \to B\). Π-types \textbf{strictly generalize} function types; they do not replace them. \texttt{Vec.replicate}'s type above literally \emph{is} \(\prod_{n : \mathtt{Nat}} \mathrm{Vec}\,\alpha\,n\), with Lean's surface syntax \texttt{(n : Nat) → Vec α n} spelling out the same thing without needing the \(\Pi\) symbol.

This is also, not by coincidence, exactly what \texttt{∀} has meant since Chapter 3. \texttt{∀ n : Nat, n ≥ 0} is a Π-type where \(B(n)\) happens to be a \emph{proposition} (\texttt{n ≥ 0 : Prop}) rather than a data type like \texttt{Vec α n} --- ``for every \texttt{n}, produce a proof of the \texttt{n}-specific statement \texttt{n ≥ 0}.'' Every \texttt{∀} written since Chapter 3 was already a dependent function, whether or not the vocabulary was available for it yet; propositions are just the special case where the family \(B\) happens to land in \texttt{Prop} instead of \texttt{Type}.

\subsection{Looking ahead}\label{sec:01-basics:03-dependent-types:7}

Chapter 11 builds a genuinely more elaborate dependent type, \texttt{Path Q : V → V → Type}, a family of types indexed by a \emph{pair} of vertices in a graph rather than by a single \texttt{Nat} --- ``the type of paths from \texttt{u} to \texttt{w},'' which differs for each choice of endpoints exactly as \texttt{Vec α n} differs for each length. Its \texttt{cons}-like constructor is a dependent function for the same reason \texttt{Vec.replicate} is one here: composing two paths is only accepted by the type-checker when their endpoints actually match, an invariant baked into the type rather than checked separately. Nothing new is needed to understand it once this section's \texttt{Fin}/\texttt{Vec} examples make sense --- it is the identical idea, with a richer index.

\begin{quote}
\textbf{Mathematical reading (optional).} For readers who already think categorically: an indexed family \texttt{B : A → Type} is exactly a functor out of the discrete category on \texttt{A}, or (thinking of \texttt{A × A}-indexed families as in the \texttt{Path} example above) an assignment of a \(\mathrm{Hom}\)-set to every pair of objects in a category --- a Π-type over such a family is a \textbf{dependent product}, a term of \(\sum_{x:A}
B(x)\) (Σ-type, next covered formally in \hyperref[sec:01-basics:05-pi-sigma-and-coc:1]{§5}) is a \textbf{dependent sum}, and both are literal categorical limits/colimits in the appropriate indexed sense --- not merely named after them by analogy.
\end{quote}

\begin{quote}
Read more: \hyperref[sec:01-basics:05-pi-sigma-and-coc:1]{§5} gives Π-types (and Σ-types) their formal typing rules, with more worked examples, rather than only the walkthrough given here.
\end{quote}

\subsection{References}\label{sec:01-basics:03-dependent-types:8}

Full citations in the \hyperref[sec:root:bibliography:1]{Bibliography}.

\begin{itemize}
\tightlist
\item
  \emph{Theorem Proving in Lean 4} (\cite{TPIL4}), ``Dependent Types'' --- the official Lean documentation's own introduction to dependent types, using this same length-indexed-vector example.
\item
  The Lean 4 source / {[}Mathlib4 API documentation{]}\cite{Mathlib4Docs} for \texttt{Fin} and \texttt{Vector} --- confirmed directly in this section via \texttt{\#print Fin} against the actual toolchain pinned in this repository's \texttt{lean\_\allowbreak{}project/\allowbreak{}lean-\allowbreak{}toolchain}.
\item
  Thompson (\cite{Thompson1991}) --- a from-scratch treatment of dependent types and Π/Σ, independent of any one proof assistant.
\item
  Chlipala (\cite{Chlipala2013}) --- Ch. 2--3 build the same length-indexed-vector idea in Coq, a useful second angle on the identical concept.
\end{itemize}
